\chapter{案例、代码与任务地图}
\label{app:case-index}

本附录按“要完成的任务”而不是按章节顺序定位材料。代码路径指向可独立运行
的程序；每个程序都把机器可读指标写入 \path{data/}，并由
\path{code/run_all.py} 统一执行。表中的“验证范围”只说明该基准实际检查了
什么，不把教学模型外推成未经计算的物理结论。

\section{十个可运行案例}

\begin{center}
\begin{tabularx}{\textwidth}{>{\raggedright\arraybackslash}p{0.07\textwidth}>{\raggedright\arraybackslash}p{0.19\textwidth}>{\raggedright\arraybackslash}p{0.30\textwidth}>{\raggedright\arraybackslash}X}
\toprule
章节 & 案例 & 程序 & 主要验证范围 \\
\midrule
5 & 谐振子 Wigner 旋转 &
\path{code/ch05_quadratic/harmonic_wigner.py} &
解析辛流、真空协方差、粒子数排序与 Monte Carlo 误差 \\
6 & Kerr 精确解--TWA &
\path{code/ch06_kerr/kerr_exact_vs_twa.py} &
短时吻合、统计误差带、独立子流与量子回复失效 \\
8 & 二聚体精确解--TWA &
\path{code/ch08_dimer/dimer_exact_vs_twa.py} &
相同初态协议、Fock 截断、轨迹守恒与短时误差 \\
11 & 压缩高斯抽样 &
\path{code/ch11_sampling/gaussian_sampling.py} &
$Q$ 压缩号约定、正常/反常矩和协方差统计检验 \\
12 & FFT 对称分裂 &
\path{code/ch12_propagation/split_step_benchmark.py} &
范数与能量、反向传播、连续步长差分的二阶收敛 \\
13 & 排序估计器 &
\path{code/ch13_observables/ordering_estimator_benchmark.py} &
相干态/热态 $n$ 与 $g^{(2)}$、jackknife 比值误差 \\
\bottomrule
\end{tabularx}
\end{center}

\clearpage
\noindent\textbf{十个可运行案例（续）}
\begin{center}
\begin{tabularx}{\textwidth}{>{\raggedright\arraybackslash}p{0.07\textwidth}>{\raggedright\arraybackslash}p{0.19\textwidth}>{\raggedright\arraybackslash}p{0.30\textwidth}>{\raggedright\arraybackslash}X}
\toprule
章节 & 案例 & 程序 & 主要验证范围 \\
\midrule
14--17 & 双分量空间场淬火 &
\path{code/ch14_17_capstone/two_component_field_quench.py} &
有限温 BdG 抽样、双场传播、相干与自旋结构因子、步长/轨迹/有限截止窗 \\
15 & BdG 淬火协方差 &
\path{code/ch15_quench/quench_covariance.py} &
稳定模式的辛归一、正常矩和反常矩 \\
16 & 集体相位扩散 &
\path{code/ch16_phase/phase_diffusion.py} &
高斯零模公式、抽样包络与解析协方差 \\
17 & 随机条纹相位标度 &
\path{code/ch17_order/stripe_scaling_toy.py} &
有序全局相位与有限相关长度畴的有限尺寸分析管线 \\
\bottomrule
\end{tabularx}
\end{center}

\section{按任务反查正文}

\begin{center}
\begin{tabularx}{\textwidth}{>{\raggedright\arraybackslash}p{0.30\textwidth}>{\raggedright\arraybackslash}X>{\raggedright\arraybackslash}X}
\toprule
任务 & 首要位置 & 必须同时核对 \\
\midrule
固定 $(x,p_x)$、$(Q,P)$ 与测度 & 第2章、附录A & 卷首符号表 \\
构造 Weyl 符号与排序修正 & 第2、4章 & 第10、13章的场论版本 \\
识别被 TWA 删除的项 & 第3、7章 & 第6、8章的精确基准 \\
判断高占据近似是否可控 & 第7、9章 & 第12章数值误差与第18章停止条件 \\
选择有限模式与真空基线 & 第10章 & 附录C的截止扫描 \\
抽样热态、压缩态或 BdG 态 & 第11、15章 & 附录B的辛协方差 \\
实现 FFT 场传播 & 第12章 & 第10章归一化、附录C测试矩阵 \\
计算密度、$g^{(2)}$ 与结构因子 & 第13章 & 第17章条纹乘积接触项 \\
区分 Rabi 与 Raman SOC & 第14章 & 第17章对称性与钉扎诊断 \\
区分去相位、预热化和热化 & 第16章 & 第9、18章的方法边界 \\
检验条纹序与超固态候选 & 第17章 & 尺寸、刚度、钉扎和独立方法 \\
选择开放系统或替代方法 & 第18章 & 附录E的分主题文献 \\
\bottomrule
\end{tabularx}
\end{center}

\section{证据文件的读取顺序}

复核一个案例时，先读对应 JSON 的 \texttt{parameters} 与
\texttt{validation.checks}，再核对图注使用的观测量和误差条，最后查看
\path{data/test_summary.json} 中记录的环境、文件哈希和总入口状态。程序零退出
只说明执行完成；阈值、实测值和通过条件均被保存，才构成结构化基准。对于
双分量空间场案例，还应区分时间步长扫描、轨迹数敏感度和三档独立重新抽样的
有限截止窗压力测试；三者回答不同问题，不能合并成一个“已收敛”的标签。
